\documentclass[11pt]{article}
\usepackage[margin=1in]{geometry}
\usepackage[T1]{fontenc}
\usepackage[utf8]{inputenc}
\IfFileExists{polski.sty}{\usepackage{polski}}{}
\usepackage{xcolor}
\usepackage{hyperref}
\hypersetup{
    colorlinks,
    citecolor=black,
    filecolor=black,
    linkcolor=black,
    urlcolor=black
}
\usepackage{amsmath, amssymb, amsthm}
\usepackage{algpseudocode}
\usepackage{algorithm}
\usepackage{indentfirst}
\usepackage{booktabs}
\usepackage{pgfplots}
\pgfplotsset{compat=1.16}
% Wiersze CSV, w ktorych kolumna #1 ma wartosc rozna od #2, sa pomijane.
\pgfplotsset{
    discard if not/.style 2 args={
        x filter/.append code={
            \edef\tempa{\thisrow{#1}}%
            \edef\tempb{#2}%
            \ifx\tempa\tempb\else\def\pgfmathresult{inf}\fi
        }
    },
    every axis/.append style={grid=major, grid style={gray!25},
        legend cell align=left, legend style={font=\scriptsize}},
    % dwa wykresy obok siebie: drugi z opcjami at=..., xshift=2cm, anchor=south west
    dwa/.style={scale only axis, width=0.36\textwidth, height=4.6cm},
    prawy/.style={at={(lewy.south east)}, xshift=2cm, anchor=south west},
}
\renewcommand{\baselinestretch}{1.3}

\makeatletter
\renewcommand{\fnum@algorithm}{Algorytm}
\makeatother

\newtheorem{theorem}{Twierdzenie}
\newtheorem{lemma}{Lemat}
\newtheorem{fact}{Fakt}
\newtheorem{corollary}{Wniosek}
\renewcommand{\proofname}{Dowód}

\algrenewcommand{\algorithmicrequire}{\textbf{Wejście:}}
\algrenewcommand{\algorithmicensure}{\textbf{Wynik:}}

\newcommand{\todo}[1]{{\color{red} \huge  \textbf{TODO: #1}}}

\newcommand{\bazaind}{\noindent \textbf{Baza indukcji. }}
\newcommand{\krokind}{\noindent \textbf{Krok indukcyjny. }}

\title{Teoria Algorytmów i Obliczeń -- projekt}
\author{Mateusz Kozłowski, Mateusz Siciarek, Mateusz Stawicki}
\date{\today}

\begin{document}

\maketitle

\tableofcontents
\newpage

\section{Wstęp}

Dokument opisuje problem podziału zbioru na \(k\) podzbiorów oraz sposoby jego rozwiązania. Propozycje rozwiązań są przedstawione w dokumencie jako opisy algorytmów wraz z dowodem ich złożoności. 

\section{Opis problemu, założenia i definicje}

\subsection{Formalny opis problemu oraz założenia}

Przyjmijmy, że wejściem problemu jest $n$-elementowy multizbiór $X = \{a_1, a_2, \dots, a_n\}$. Zakładamy, że wszystkie elementy są nieujemnymi liczbami całkowitymi, czyli:
\begin{equation}
    \forall_{i \in \{1, \dots, n\}} : a_i \in \mathbb{N}_0,
    \label{eq:positive_elems}
\end{equation}
gdzie $\mathbb{N}_0$ oznacza zbiór liczb naturalnych wraz z zerem. Zwróćmy uwagę, że elementy multizbioru rozróżniamy poprzez ich indeksy - dwie jednakowe wartości liczbowe stanowią osobne elementy.

Oznaczmy przez $[n] = \{1, \dots, n\}$ zbiór indeksów elementów multizbioru $X$. Całkowitą sumę wszystkich elementów oznaczamy jako:
\begin{equation}
    S_X = \sum_{i=1}^n a_i.
\end{equation}

Dla dowolnego podzbioru indeksów $I \subseteq [n]$, sumę odpowiadających mu elementów definiujemy jako:
\begin{equation}
    S_I = \sum_{i \in I} a_i.
\end{equation}

\subsubsection{Opis problemu - podział na $k$ niepustych podzbiorów}

Rozważmy problem polegający na podziale multizbioru $X$ na $k \ge 2$ podzbiorów. Poszukujemy partycji zbioru indeksów $[n]$ na $k$ rozłącznych i niepustych podzbiorów $(A_1, A_2, \dots, A_k)$ spełniających warunki:
\begin{equation}
    \bigcup_{j=1}^k A_j = [n] \quad \text{oraz} \quad A_i \cap A_j = \emptyset \quad \text{dla wszystkich } i \neq j.
\end{equation}

Chcemy zminimalizować największą różnicę między sumami elementów w dowolnych dwóch podzbiorach partycji. Dla dwóch sum definiujemy ich różnicę bezwzględną jako:
\begin{equation}
    \Delta(S_I, S_J) = |S_I - S_J|.
\end{equation}
Funkcję celu definiujemy wtedy jako:
\begin{equation}
    \Delta(A_1, \dots, A_k) = \max_{1 \le i < j \le k} \Delta(S_{A_i}, S_{A_j}).
\end{equation}

Problem optymalizacyjny sprowadza się do znalezienia podziału $(A_1, \dots, A_k)$ minimalizującego wartość tak zdefiniowanej funkcji $\Delta$:
\begin{equation}
    \min_{(A_1, \dots, A_k)} \Delta(A_1, \dots, A_k).
\end{equation}

W tym przypadku \textbf{podział idealny} zachodzi wtedy i tylko wtedy, gdy suma elementów w każdym z $k$ podzbiorów jest identyczna i wynosi:
\begin{equation}
    S_{A_1} = S_{A_2} = \dots = S_{A_k} = \frac{1}{k} S_X,
\end{equation}
co odpowiada wartości funkcji celu $\Delta(A_1, \dots, A_k) = 0$. Warunkiem koniecznym (choć niewystarczającym) do istnienia podziału idealnego jest podzielność całkowitej sumy $S_X$ przez liczbę podzbiorów $k$, czyli $S_X \bmod k = 0$.

\section{Algorytmy -- opisy i dowody}

\subsection{Algorytm dokładny przeglądu zupełnego}

\subsubsection{Specyfikacja problemu}

Wejściem jest ciąg \(a_1,\ldots,a_n\) spełniający założenie \eqref{eq:positive_elems} oraz liczba części \(k\geq 2\).
Poszczególne wystąpienia wartości rozróżniamy za pomocą ich indeksów.

Algorytm ma ustalić, czy istnieje \textbf{podział idealny} zdefiniowany w sekcji problemowej, czyli partycja \(\mathcal{P}=(A_1,\ldots,A_k)\) zbioru indeksów \([n]\) na \(k\) niepustych, parami rozłącznych części, której sumy części spełniają \(S_{A_1}=S_{A_2}=\cdots=S_{A_k}=S_X/k\), co odpowiada wartości zero funkcji celu \(\Delta\). Jeżeli taki podział istnieje, algorytm zwraca jedną partycję go realizującą. W przeciwnym przypadku zwraca komunikat \texttt{BRAK}; w szczególności, jeśli \(n<k\), dopuszczalny podział nie istnieje.

Podział idealny, o ile istnieje, nie musi być jednoznaczny -- różne rozkłady indeksów na części mogą dawać równe sumy, a ta sama partycja ma \(k!\) różniących się jedynie numeracją części reprezentacji. Algorytm zwraca \textbf{pierwsze napotkane} rozwiązanie w kolejności generowania kandydatów i natychmiast kończy działanie; pozostałych kandydatów nie bada już wcale. Odpowiedź \texttt{BRAK} oznacza, że żaden z wszystkich kandydatów nie reprezentuje podziału idealnego.

\subsubsection{Idea algorytmu i generowanie kandydatów}

Przypomnijmy, że \(S_X=\sum_{i=1}^{n}a_i\) oznacza sumę wszystkich elementów wejścia. Jeżeli \(n<k\), nie istnieje podział \([n]\) na \(k\) niepustych części. Jeżeli \(S_X\) nie jest podzielna przez \(k\), podział idealny również nie istnieje, ponieważ z równości sum wszystkich części wynika \(S_X=k\cdot S_{A_1}\).

W pozostałych przypadkach wystarczy znaleźć podział \([n]\) na \(k\) niepustych części o sumach równych \(\mathrm{cel}=S_X/k\). Podział reprezentujemy wektorem \(b=(b_1,\ldots,b_n)\in[k]^n\), gdzie \(b_i=j\) oznacza \(i\in A_j\). Rodzina zbiorów
\begin{equation}
A_j=\{i\in[n]:\ b_i=j\}, \qquad j\in[k],
\end{equation}
jest wówczas zawsze podziałem zbioru \([n]\) na \(k\) rozłącznych części -- zmienną pozostaje jedynie niepustość tych części.

Dla \(n\geq k\) definiujemy zbiór kandydatów:
\begin{equation}
\mathcal{V}_{n,k}=[k]^n.
\end{equation}
Kandydat \(b\) jest \textbf{dopuszczalny}, jeśli każda z części \(A_1,\ldots,A_k\) jest niepusta, czyli jeśli każda wartość z \([k]\) występuje w \(b\) co najmniej raz. Niepustość części testujemy za pomocą liczników wystąpień wartości w wektorze \(b\), a nie za pomocą sum części. Jest to istotne, ponieważ elementy wejścia mogą być zerami: dla wejścia złożonego z samych zer zachodzi \(S_X=0\) oraz \(\mathrm{cel}=0\), więc sam test sum nie odróżniłby pustej części od części złożonej z elementów zerowych.

Liczba kandydatów wynosi:
\begin{equation}
m=|\mathcal{V}_{n,k}|=k^n.
\end{equation}
Instrukcja \texttt{for each b in [k]\^{}n} oznacza generowanie kolejno wszystkich wektorów z \([k]^n\), każdego dokładnie raz, bez przechowywania całego zbioru. Można to zrealizować przez zliczanie \(k\)-czne na \(n\) pozycjach: zaczynamy od wektora \((1,\ldots,1)\), a następnie zwiększamy licznik o jeden, kończąc po przetworzeniu układu \((k,\ldots,k)\). Pojedyncze zwiększenie wymaga najwyżej \(O(n)\) zmian pozycji, jednak łączna liczba zmian pozycji podczas całego zliczania wynosi
\begin{equation}
k^n+k^{n-1}+\cdots+k+1=\frac{k^{n+1}-1}{k-1}<2k^n,
\end{equation}
zatem koszt zamortyzowany generowania jednego kandydata wynosi \(O(1)\).

Kolejność generowania jest deterministyczna, więc pojęcie \textbf{pierwszego napotkanego} rozwiązania jest dobrze określone: dla ustalonego wejścia algorytm zwraca zawsze ten sam wynik, nawet jeśli rozwiązań istnieje wiele. Zauważmy wreszcie, że jedna partycja \((A_1,\ldots,A_k)\) ma \(k!\) reprezentacji wektorowych w \([k]^n\) (wynikających z permutacji numerów części). Redundancja ta nie wpływa na poprawność -- wszystkie reprezentacje tej samej partycji dają identyczne rodziny sum -- ani na rząd złożoności, ponieważ algorytm przegląda wszystkie \(k^n\) wektorów z \(\mathcal{V}_{n,k}\) niezależnie od tego, ile z nich reprezentuje ten sam podział.

\subsubsection{Pseudokod}

\begin{verbatim}
BRUTE-K-PODZIAL(a[1..n], k):
    if n < k:
        return BRAK

    S_X := 0
    for i := 1 to n:
        S_X := S_X + a[i]

    if S_X mod k != 0:
        return BRAK

    cel := S_X / k

    for each b in [k]^n:
        for j := 1 to k:
            S[j] := 0
            licz[j] := 0

        for i := 1 to n:
            S[b[i]] := S[b[i]] + a[i]
            licz[b[i]] := licz[b[i]] + 1

        dopuszczalny := true
        for j := 1 to k:
            if licz[j] = 0:
                dopuszczalny := false

        if dopuszczalny:
            idealny := true
            for j := 1 to k:
                if S[j] != cel:
                    idealny := false

            if idealny:
                // pierwsze trafienie: wynik zwracamy natychmiast,
                // pozostalych kandydatow nie badamy
                for j := 1 to k:
                    A[j] := {i in [n] : b[i] = j}
                return (A[1], ..., A[k])

    return BRAK
\end{verbatim}

Dzielenie \(S_X/k\) daje tutaj wynik całkowity, ponieważ wartości \(S_X\) niedzielne przez \(k\) zostały wcześniej wykluczone.

\subsubsection{Dowód częściowej poprawności}

Pokażemy, że każdy wynik zwrócony przez algorytm jest zgodny ze specyfikacją. Niezmienniki pętli udowodnimy indukcyjnie względem liczby wykonanych iteracji.

\paragraph{Poprawność początkowego odrzucenia.}

Jeżeli \(n<k\), nie można podzielić \(n\)-elementowego zbioru \([n]\) na \(k\) niepustych części. Ponadto istnienie podziału idealnego pociąga za sobą:
\begin{equation}
S_X=\sum_{j=1}^{k}S_{A_j}=k\cdot S_{A_1}.
\end{equation}
Wobec całkowitości \(S_{A_1}\) liczba \(S_X\) musi być podzielna przez \(k\). Oba początkowe warunki odrzucenia są zatem poprawne.

\paragraph{Niezmiennik pętli obliczającej sumę \(S_X\).}

Niech \(\Sigma_j\) oznacza wartość zmiennej \texttt{S\_X} po \(j\) iteracjach pierwszej pętli. Wykażemy indukcyjnie, że:
\begin{equation}
\Sigma_j=\sum_{i=1}^{j}a_i, \qquad 0\leq j\leq n.
\end{equation}
\bazaind Dla \(j=0\) zmienna \texttt{S\_X} ma wartość \(0\), równą sumie pustej.

\krokind Załóżmy, że równość zachodzi dla pewnego \(j<n\). W następnej iteracji do zmiennej dodawane jest \(a_{j+1}\), więc:
\begin{equation}
\Sigma_{j+1}=\Sigma_j+a_{j+1}=\sum_{i=1}^{j+1}a_i.
\end{equation}
Na mocy indukcji po zakończeniu pętli \(\Sigma_n\) jest sumą wszystkich elementów wejścia, czyli \(\Sigma_n=S_X\).

\paragraph{Niezmiennik pętli wewnętrznej.}

Ustalmy kandydata \(b\in[k]^n\). Niech \(S_j^{(t)}\) oraz \(\mathrm{licz}_j^{(t)}\) oznaczają wartości zmiennych \texttt{S[j]} i \texttt{licz[j]} po przetworzeniu pierwszych \(t\) pozycji. Wykażemy indukcyjnie, że dla każdego \(j\in[k]\) jednocześnie:
\begin{equation}
S_j^{(t)}=\sum_{\substack{1\leq i\leq t\\b_i=j}}a_i, \qquad \mathrm{licz}_j^{(t)}=\bigl|\{i\leq t:\ b_i=j\}\bigr|, \qquad 0\leq t\leq n.
\end{equation}
\bazaind Dla \(t=0\) wszystkie sumy i liczniki mają wartość \(0\), równą sumie pustej i mocy zbioru pustego.

\krokind Załóżmy, że niezmiennik zachodzi dla pewnego \(t<n\). Oznaczmy \(j^{*}=b_{t+1}\). W następnej iteracji do \texttt{S[j*]} dodawane jest \(a_{t+1}\), a \texttt{licz[j*]} zwiększane o jeden, natomiast zmienne o indeksach innych niż \(j^{*}\) pozostają bez zmian. Indeks \(t+1\) należy właśnie do zbioru \(\{i\leq t+1:\ b_i=j^{*}\}\) i do żadnego innego zbioru tej postaci, więc:
\begin{equation}
S_{j^{*}}^{(t+1)}=S_{j^{*}}^{(t)}+a_{t+1}=\sum_{\substack{1\leq i\leq t+1\\b_i=j^{*}}}a_i, \qquad \mathrm{licz}_{j^{*}}^{(t+1)}=\bigl|\{i\leq t+1:\ b_i=j^{*}\}\bigr|,
\end{equation}
a dla \(j\neq j^{*}\) obie równości przenoszą się z założenia indukcyjnego. Na mocy indukcji po zakończeniu pętli, dla każdego \(j\in[k]\):
\begin{equation}
\texttt{S[j]}=S_j^{(n)}=\sum_{i\in A_j}a_i=S_{A_j}, \qquad \texttt{licz[j]}=\mathrm{licz}_j^{(n)}=|A_j|,
\end{equation}
gdzie \(A_j=\{i\in[n]:\ b_i=j\}\).

\paragraph{Poprawność testu dopuszczalności.}

Z udowodnionego niezmiennika test \texttt{licz[j] = 0} odrzuca kandydata wtedy i tylko wtedy, gdy któraś z części \(A_1,\ldots,A_k\) jest pusta. Kandydat przechodzi test dokładnie wtedy, gdy jest surjekcją z \([n]\) na \([k]\).

Zwróćmy uwagę, dlaczego liczniki są konieczne. Dla wejścia złożonego z samych zer mamy \(S_X=0\) oraz \(\mathrm{cel}=0\). Wektor \(b=(1,\ldots,1)\) spełnia wówczas warunek \(\texttt{S[j]}=\mathrm{cel}\) dla części \(A_1=[n],\ A_2=\cdots=A_k=\emptyset\), choć nie jest dopuszczalny. Sam test sum nie odróżniłby więc pustej części od części złożonej z elementów zerowych i mógłby zwrócić podział naruszający warunek niepustości.

\paragraph{Poprawność zwróconego podziału.}

Załóżmy, że algorytm zwraca ciąg \((A_1,\ldots,A_k)\). Z konstrukcji \(A_j=\{i\in[n]:\ b_i=j\}\) każda pozycja trafia do dokładnie jednej części, zatem:
\begin{equation}
\bigcup_{j=1}^{k}A_j=[n], \qquad A_i\cap A_j=\emptyset\quad(i\neq j).
\end{equation}
Ponieważ kandydat przeszedł test dopuszczalności, wszystkie części są niepuste. Z niezmiennika pętli wewnętrznej oraz testu \texttt{S[j] = cel} otrzymujemy ponadto:
\begin{equation}
S_{A_j}=\frac{S_X}{k} \qquad\text{dla każdego }j\in[k].
\end{equation}
Zwrócony wynik jest więc dopuszczalnym podziałem idealnym.

\paragraph{Własność zwróconego pierwszego rozwiązania.}

Zwrócony podział jest dopuszczalny i idealny niezależnie od tego, ile rozwiązań istnieje. Jeżeli rozwiązań jest wiele, algorytm zwraca to z nich, którego reprezentacja wektorowa pojawia się najwcześniej w kolejności zliczania \(k\)-cznego. Reprezentacje różniące się jedynie permutacją numerów części tej samej partycji dają partycje o identycznych rodzinach sum, więc wybór konkretnej reprezentacji nie wpływa na zgodność wyniku ze specyfikacją. Wcześniejszy zwrot pomija jedynie kandydatów, którzy -- na mocy niezmiennika pętli zewnętrznej -- i tak nie są rozwiązaniami.

\paragraph{Niezmiennik pętli zewnętrznej.}

Niech \(b^{(1)},\ldots,b^{(m)}\), \(m=k^n\), oznaczają wektory z \([k]^n\) w kolejności ich generowania. Wykażemy indukcyjnie względem \(t\), że po przetworzeniu pierwszych \(t\) kandydatów, jeżeli algorytm nie zwrócił wyniku i kontynuuje działanie, zachodzi:
\begin{equation}
\forall r\in\{1,\ldots,t\}\qquad b^{(r)}\ \text{nie jest dopuszczalnym wektorem podziału idealnego},
\end{equation}
czyli \(b^{(r)}\) nie spełnia łącznie warunku suriektywności oraz równości wszystkich sum części i wartości \(\mathrm{cel}\).

\noindent \textbf{Baza indukcji.} Dla \(t=0\) nie sprawdzono jeszcze żadnego kandydata. Warunek zachodzi, ponieważ kwantyfikacja dotyczy zbioru pustego.

\noindent \textbf{Krok indukcyjny.} Załóżmy, że niezmiennik zachodzi po \(t<m\) iteracjach i algorytm kontynuuje działanie. W następnej iteracji przetwarzany jest wektor \(b^{(t+1)}\). Na podstawie udowodnionych niezmienników algorytm poprawnie wyznacza sumy i liczniki wszystkich części, a więc poprawnie rozstrzyga dopuszczalność i idealność kandydata.

Jeżeli \(b^{(t+1)}\) jest dopuszczalnym wektorem podziału idealnego, algorytm zwraca poprawny podział i kończy działanie -- niezmiennik nie musi już obowiązywać. Jeżeli nie jest, algorytm kontynuuje działanie i własność ta zachodzi dla \(b^{(t+1)}\). Łącząc ją z założeniem indukcyjnym, otrzymujemy niezmiennik dla \(t+1\).

Z niezmiennika wynika w szczególności, że jeżeli algorytm zwraca wynik przy kandydacie \(b^{(t+1)}\), to \(t+1\) jest \textbf{najmniejszym} indeksem o tej własności -- zwrócony wektor jest istotnie pierwszym trafieniem w kolejności generowania.

\paragraph{Poprawność końcowego wyniku \texttt{BRAK}.}

Jeżeli algorytm dociera do ostatniej instrukcji \texttt{return BRAK}, przetworzył wszystkie \(m=k^n\) wektorów. Z niezmiennika pętli zewnętrznej wynika, że żaden z nich nie jest dopuszczalnym wektorem podziału idealnego.

Załóżmy przeciwnie, że istnieje dopuszczalny podział idealny \((A_1^{*},\ldots,A_k^{*})\). Rozważmy wektor \(b^{*}\) dany przez \(b^{*}_i=j\iff i\in A_j^{*}\). Ponieważ wszystkie części są niepuste, \(b^{*}\) jest surjekcją na \([k]\), a zatem \(b^{*}\in[k]^n\) i \(b^{*}=b^{(r)}\) dla pewnego \(r\in\{1,\ldots,m\}\). Dla tego kandydata wszystkie liczniki są dodatnie, a wszystkie sumy części równe \(S_X/k=\mathrm{cel}\), więc algorytm zwróciłby wynik najpóźniej w \(r\)-tej iteracji, co przeczy dojściu do instrukcji \texttt{return BRAK}.

Końcowy wynik \texttt{BRAK} jest więc poprawny. Łącznie z poprzednimi przypadkami dowodzi to częściowej poprawności algorytmu.

\subsubsection{Dowód własności stopu}

Dla \(n<k\) algorytm zatrzymuje się natychmiast. Dla \(n\geq k\) pierwsza pętla ma dokładnie \(n\) iteracji. Jeśli \(S_X\) nie dzieli się przez \(k\), algorytm następnie się zatrzymuje.

W pozostałym przypadku generator produkuje dokładnie \(m=k^n\) kandydatów, każdego w skończonym czasie -- każde zwiększenie licznika \(k\)-cznego wymaga najwyżej \(n\) zmian pozycji. Po przetworzeniu \(t\) kandydatów liczba pozostałych wynosi \(m-t\); przy kontynuowaniu pętli maleje ona o jeden i jest nieujemna. Dla każdego kandydata pętla wyznaczająca sumy i liczniki wykonuje dokładnie \(n\) iteracji, natomiast inicjalizacja tablic oraz testy wykonują łącznie co najwyżej \(3k\) iteracji pętli po \(k\) kroków. Wcześniejszy zwrot wyniku może jedynie skrócić działanie algorytmu i nie wpływa na ten rachunek.

Wszystkie pętle są zatem skończone. Utworzenie zbiorów wynikowych wymaga przejrzenia \(n\) pozycji, więc także kończy się w skończonym czasie. Algorytm zatrzymuje się dla każdego wejścia. W połączeniu z częściową poprawnością otrzymujemy całkowitą poprawność.

\subsubsection{Złożoność czasowa}

\paragraph{Model jednostkowego kosztu.}

Najpierw prowadzimy analizę w modelu jednostkowego kosztu operacji arytmetycznych, porównań i dostępu do elementów tablicy, zakładając, że wartości wejściowe, ich sumy oraz \(k\) mieszczą się w słowie pamięci.

Obliczenie \(S_X\) wymaga \(\Theta(n)\) czasu. Kandydatów jest \(m=k^n\), a koszt zamortyzowany wygenerowania jednego z nich wynosi \(O(1)\). Sprawdzenie jednego kandydata wymaga \(\Theta(n)\) czasu na wyznaczenie sum i liczników oraz \(\Theta(k)\) czasu na inicjalizację tablic \texttt{S} i \texttt{licz} oraz na test dopuszczalności, który przegląda wszystkie pozycje liczników zawsze w całości, niezależnie od ich wartości; test idealności wykonywany jedynie dla kandydatów dopuszczalnych mieści się w tej samej granicy. Utworzenie zbiorów wynikowych wymaga \(O(n)\) czasu i następuje najwyżej raz. Stąd:
\begin{equation}
T_{\mathrm{pes}}(n,k)=O\bigl(n+(n+k)\,k^n\bigr).
\end{equation}
Aby wykazać dolne oszacowanie, dla dowolnych \(n\geq k\geq2\) rozważmy wejście:
\begin{equation}
a_1=\cdots=a_{n-1}=k, \qquad a_n=k(n+1).
\end{equation}
Wtedy:
\begin{equation}
S_X=k(n-1)+k(n+1)=2kn, \qquad \frac{S_X}{k}=2n,
\end{equation}
zatem \(S_X\) jest podzielna przez \(k\) i test wstępny nie odrzuca tego wejścia. Jednak część zawierająca indeks \(n\) ma wartość co najmniej \(k(n+1)>2n\), więc jej suma różni się od \(\mathrm{cel}=2n\) i podział idealny nie istnieje. Wczesne zwrócenie wyniku jest niemożliwe, a algorytm sprawdza wszystkich \(m=k^n\) kandydatów, wydając na każdym kandydacie \(\Theta(n+k)\) czasu. Otrzymujemy:
\begin{equation}
T_{\mathrm{pes}}(n,k)=\Omega\bigl((n+k)\,k^n\bigr).
\end{equation}
Łącząc oszacowania oraz korzystając z \(k\leq n\), mamy ostatecznie:
\begin{equation}
\boxed{T_{\mathrm{pes}}(n,k)=\Theta\bigl((n+k)\,k^n\bigr)=\Theta\bigl(n\,k^n\bigr)\quad\text{dla }2\leq k\leq n.}
\end{equation}
Zwróćmy uwagę, że oszacowanie dotyczy najgorszego przypadku: wejść bez rozwiązania albo takich, których rozwiązanie pojawia się późno w kolejności generowania. Wejścia z rozwiązaniem wcześniejszym w tej kolejności mogą kończyć się szybciej, algorytm nie ma jednak gwarancji znalezienia rozwiązania przed pełnym przeglądem.

\paragraph{Kryterium logarytmiczne.}

W modelu jednostkowym założyliśmy, że dodawanie i porównywanie liczb ma koszt stały. W kryterium logarytmicznym operacje te \textbf{nie są atomowe}: ich koszt rośnie liniowo wraz z długością binarnego zapisu operandów. Przyjmijmy, zgodnie z konwencją przyjętą dla algorytmu aproksymacyjnego, że
\begin{equation}
b_i=\max\bigl(1,\lceil\log_2(a_i+1)\rceil\bigr), \qquad b=\max_{1\leq i\leq n}b_i, \qquad s=b+\lceil\log_2 n\rceil,
\end{equation}
gdzie \(b_i\) jest liczbą bitów zapisu \(a_i\), a \(s\) górnym ograniczeniem długości zapisu sum części, wartości \(S_X\) oraz \(\mathrm{cel}\); istotnie, \(S_j\leq S_X<n2^b\). Indeksy, liczniki wystąpień (\(\mathrm{licz}[j]\leq n\)) oraz \(k\) traktujemy jako mieszczące się w słowie maszynowym o koszcie \(O(1)\).

Pojedyncze dodanie \(a_{i}\) do sumy \texttt{S[j]}, porównanie \texttt{S[j]} z \(\mathrm{cel}\), test \(S_X\bmod k\) oraz obliczenie \(\mathrm{cel}=S_X/k\) operują na liczbach długości \(\Theta(s)\) bitów. Analogicznie jak w analizie algorytmu aproksymacyjnego, przyjmujemy najgorszy przypadek pojedynczej operacji danego typu: każda z tych operacji kosztuje wtedy dokładnie \(\Theta(s)\) operacji bitowych, niezależnie od konkretnej pary operandów. Zatem obliczenie \(S_X\) wymaga \(\Theta(ns)\) czasu, a sprawdzenie jednego kandydata wymaga \(n\) dodawań oraz \(O(k)\) porównań, co daje \(\Theta((n+k)\,s)\). Dla rodziny trudnej z poprzedniego akapitu algorytm wykonuje pełny przegląd wszystkich \(k^n\) kandydatów, stąd dolne i górne oszacowanie zbiegają się:
\begin{equation}
\boxed{T_{\mathrm{bit}}(n,k,b)=\Theta\bigl((n+k)\,k^n\,\bigl(b+\lceil\log_2 n\rceil\bigr)\bigr).}
\end{equation}

Zwróćmy uwagę na relację tego wyniku do rozmiaru wejścia. Rozmiar wejścia wynosi \(N=O(B)\) bitów, gdzie \(B=\sum_i b_i\), oraz \(B\geq n\). Dla rodzin wejść o wartościach ze zbioru \(\{0,1\}\) zachodzi \(b_i=1\) dla każdego \(i\), zatem \(B=\Theta(n)\) oraz \(N=\Theta(n)\). Ponieważ \(k\geq2\), czas działania algorytmu na takich wejściach spełnia
\begin{equation}
n\,k^n\geq n\,2^n=2^{\Omega(N)},
\end{equation}
zatem czas działania jest \textbf{eksponencjalny względem rozmiaru wejścia} w kryterium logarytmicznym -- już dla pesymistycznych wejść o jednubitowych wartościach. Model jednostkowy ukrywa czynnik \(s\), który w modelu logarytmicznym uwidacznia nieatomowość operacji na wielkich liczbach.
\subsection{Autorski algorytm aproksymacyjny}

\subsubsection{Specyfikacja problemu i rozmiar danych wejściowych}

Dane są liczby $a_1,\ldots,a_n$ spełniające założenie \eqref{eq:positive_elems} oraz liczba części $k\geq 2$. Wystąpienia jednakowych wartości rozróżniamy za pomocą ich oryginalnych indeksów. 

Szukamy partycji $\mathcal{P}=(A_1,\ldots,A_k)$ zbioru indeksów $[n]$ na $k$ niepustych, parami rozłącznych części, określonej w sekcji problemowej. Dla $S_j=\sum_{i\in A_j}a_i$ minimalizujemy rozpiętość sum
\begin{equation}
  \Delta(\mathcal{P})=\max_{1\leq j\leq k}S_j-\min_{1\leq j\leq k}S_j.
\end{equation}
Wielkość ta równa się $\max_{i\neq j}|S_i-S_j|$, czyli największej różnicy sum dwóch części, zgodnie ze sformułowaniem funkcji celu z sekcji problemowej. Jeśli $n<k$, dopuszczalny podział nie istnieje i algorytm zwraca komunikat \texttt{BRAK}; dla $n\geq k$ ma on zwrócić dowolny dopuszczalny podział. Przez $\Delta^*$ oznaczamy najmniejszą rozpiętość dopuszczalnego podziału dla $n\geq k$.

Rozmiar wejścia zależy zarówno od liczby elementów $n$, jak i od długości zapisu samych liczb. Niech
\begin{equation}
  b_i=\max\bigl(1,\lceil\log_2(a_i+1)\rceil\bigr),\qquad
  B=\sum_{i=1}^{n}b_i,\qquad
  b=\max_{1\leq i\leq n}b_i.
\end{equation}
Wielkość $b_i$ oznacza liczbę bitów potrzebnych do zapisania $a_i$, $B$ jest łączną długością zapisów wszystkich wartości wejściowych, a $b$ jest długością najdłuższej z nich. Każda liczba zajmuje przynajmniej jeden bit, zatem $B\geq n$. Przypadek $k>n$ kończy się trywialnie odpowiedzią \texttt{BRAK}, zatem dalej zakładamy $2\leq k\leq n$; narzut separatorów oraz zapisu parametrów $n,k$ mieści się wtedy w $O(B)$ bitów. Przyjmujemy więc, że rozmiar całego wejścia wynosi
\begin{equation}
  N=O(B)\text{ bitów}.
\end{equation}

Zakładamy, że parametry $n,k$ oraz indeksy i liczniki mieszczą się w słowie maszynowym, a operacje na nich mają koszt $O(1)$. Wartości $a_i$ oraz sumy $S_j$ przechowujemy jako liczby całkowite o dowolnej precyzji. Ponieważ
\begin{equation}
  S_j\leq\sum_{i=1}^{n}a_i<n2^b,
\end{equation}
długość zapisu dowolnej sumy nie przekracza
\begin{equation}
  s=b+\lceil\log_2 n\rceil
\end{equation}
bitów. Przyjmujemy konwencję najgorszego przypadku pojedynczej operacji danego typu: porównanie dwóch wartości wejściowych kosztuje co najwyżej $O(b)$ operacji bitowych, natomiast dodawanie lub porównywanie sum co najwyżej $O(s)$ operacji bitowych. Koszt tych działań zależy więc liniowo od długości zapisu liczb w bitach, a tym samym logarytmicznie od ich wartości.


\subsubsection{Idea algorytmu}

Sortujemy pary $(a_i,i)$ malejąco według wartości, a przy równych
wartościach rosnąco według oryginalnego indeksu. Pierwsze $k$ par
umieszczamy po jednej w częściach $A_1,\ldots,A_k$. Dzięki temu każda
część jest niepusta już po inicjalizacji.

Pozostałe pary przetwarzamy w kolejnych blokach po $k$ elementów,
ostatni blok może być krótszy. W obrębie każdego bloku przechodzimy
od największej wartości do najmniejszej. Dla każdej pary wybieramy część
o najmniejszej aktualnej sumie, dopisujemy do niej oryginalny indeks
i aktualizujemy sumę. Wybór oraz aktualizację minimum obsługuje kopiec
minimum przechowujący pary $(S_j,j)$. Przy równych sumach wybieramy
najmniejszy numer części. Ta sama część może otrzymać kilka elementów
jednego bloku, a inna żadnego.

\subsubsection{Pseudokod}

Po posortowaniu $v_p$ oznacza wartość, a $\operatorname{idx}_p$ oryginalny
indeks elementu na pozycji $p$; pary spełniają
$v_1\geq v_2\geq\cdots\geq v_n$ i
$v_p=a_{\operatorname{idx}_p}$. Tablica $\operatorname{idx}$ jest
permutacją $[n]$.

\begin{verbatim}
MALEJACY-PRZYDZIAL(a[1..n], k):
    if n < k:
        return BRAK 

    (v, idx) := sortuj pary (a[i], i) malejąco według wartości, 
                a przy remisie rosnąco według indeksu

    // Przydział po jednym elemencie zapewnia niepustość każdej części
    for j := 1 to k: 
        A[j] := lista zawierajaca idx[j] 
        S[j] := v[j]

    H := BUILD-MIN-HEAP({(S[j], j) : j = 1,...,k})
    p := k + 1 // Pozycja pierwszego elementu, który nie został jeszcze przypisany

    while p <= n: 
        // Koniec bloku liczącego najwyżej k elementów - ostatni blok może być krótszy
        koniec := min(n, p + k - 1) 
       
        for q := p to koniec: // przetwarzanie bloku
            (suma, j) := EXTRACT-MIN(H)
            dopisz idx[q] do A[j]
            S[j] := suma + v[q]
            INSERT(H, (S[j], j))

        p := koniec + 1

    return (A[1], ..., A[k])
\end{verbatim}

\subsubsection{Dowód poprawności i własność stopu}

Jeśli $n<k$, podział $n$ indeksów na $k$ niepustych części nie istnieje,
więc odpowiedź \texttt{BRAK} jest poprawna. Dalej zakładamy $n\geq k$.
Po inicjalizacji każda część zawiera dokładnie jeden inny indeks, sumy
odpowiadają wartościom przypisanych elementów, a kopiec zawiera dokładnie
jedną parę $(S_j,j)$ dla każdej części.

Przed każdym następnym przydziałem przyjmujemy następujący niezmiennik:
indeksy z pozycji $1,\ldots,q-1$ posortowanej tablicy zostały przypisane
dokładnie raz, indeksy z pozycji $q,\ldots,n$ jeszcze nie zostały
przypisane, każda część jest niepusta, zapisane sumy są poprawne,
a kopiec zawiera po jednej aktualnej parze $(S_j,j)$ dla każdej części.
Niezmiennik zachodzi przed pierwszym przydziałem, gdy $q=k+1$.

Operacja \texttt{EXTRACT-MIN} usuwa parę części o najmniejszej sumie.
Dopisanie $\operatorname{idx}_q$ przypisuje dokładnie jeden nowy indeks,
a $S_j:=S_j+v_q$ daje sumę elementów w tej części. Operacja
\texttt{INSERT} przywraca jej jedyną aktualną parę w kopcu. Pozostałe
części pozostają bez zmian, zatem niezmiennik zachodzi przed kolejnym
przydziałem. Granice bloków jedynie dzielą tę samą kolejność pozycji
$k+1,\ldots,n$; żadna pozycja nie jest pomijana ani przetwarzana
dwukrotnie. Po przetworzeniu ostatniej pozycji wszystkie indeksy należą
do dokładnie jednej części, części są niepuste, a ich sumy są poprawne.
Otrzymany podział jest więc dopuszczalny.

Każdy obieg pętli zewnętrznej zwiększa $p$ o co najmniej $1$,
a każda pętla wewnętrzna ma co najwyżej $k$ przebiegów. Sortowanie
i operacje na skończonym kopcu również kończą się, więc algorytm się
zatrzymuje.

\subsubsection{Gwarancja jakości}
Niech $a_{\max}=\max_i a_i$. Dla otrzymanego podziału zachodzi
\begin{equation}
  \boxed{\Delta_{\mathrm{alg}}\leq a_{\max}.}
\end{equation}

Po początkowym przydziale każda część zawiera jeden element wejściowy.
Ponieważ wszystkie wartości należą do przedziału $[0,a_{\max}]$,
początkowa rozpiętość nie przekracza $a_{\max}$.

Rozważmy następny przydział. Niech $m$ i $M$ oznaczają odpowiednio
najmniejszą i największą sumę przed przydziałem, a $m'$ i $M'$
odpowiednie wartości po przydziale. Dotychczasowa rozpiętość wynosi
$\Delta=M-m$.

Dodajemy element o wartości $x\geq0$ do części o najmniejszej sumie $m$.
Żadna suma nie maleje, zatem
\begin{equation}
  m'\geq m.
\end{equation}
Sumy niezmienionych części nie przekraczają $M$, natomiast suma
zmienionej części wynosi $m+x$. Dlatego
\begin{equation}
  M'\leq\max(M,m+x).
\end{equation}

Korzystając z tych dwóch nierówności, otrzymujemy
\begin{equation}
  \begin{aligned}
    \Delta'=M'-m'
    &\leq M'-m\\
    &\leq \max(M,m+x)-m\\
    &=\max(M-m,m+x-m)\\
    &=\max(M-m,x)\\
    &=\max(\Delta,x).
  \end{aligned}
\end{equation}
Pierwsza nierówność wynika z zastąpienia odejmowanej wartości $m'$
wartością $m\leq m'$, co może zwiększyć różnicę.
W drugiej nierówności zastępujemy $M'$ jego górnym ograniczeniem
$\max(M,m+x)$.
Następnie odejmujemy $m$ od obu argumentów maksimum.

Jeżeli przed przydziałem zachodzi $\Delta\leq a_{\max}$,
to wobec $x\leq a_{\max}$ otrzymujemy
\begin{equation}
  \Delta'\leq\max(\Delta,x)\leq a_{\max}.
\end{equation}
Ponieważ nierówność zachodzi po początkowym przydziale,
indukcja po kolejnych przydziałach daje
$\Delta_{\mathrm{alg}}\leq a_{\max}$.

Wynik algorytmu jest dopuszczalnym podziałem, więc
$\Delta^*\leq\Delta_{\mathrm{alg}}$.
Ponadto $\Delta^*\geq0$, zatem otrzymujemy addytywną gwarancję
\begin{equation}
  0\leq\Delta_{\mathrm{alg}}-\Delta^*
  \leq\Delta_{\mathrm{alg}}
  \leq a_{\max}.
\end{equation}

Nie istnieje natomiast stały multiplikatywny współczynnik aproksymacji
rozpiętości. Dla $t\geq1$ weźmy $k=3$ i liczby posortowane malejąco:
\begin{equation}
  6t+1,\ 3t,\ 3t,\ 2t,\ 2t,\ 2t.
\end{equation}
Początkowe sumy wynoszą $(6t+1,3t,3t)$. Trzy kolejne przydziały
dają odpowiednio $(6t+1,5t,3t)$, $(6t+1,5t,5t)$ i
$(6t+1,7t,5t)$, zatem $\Delta_{\mathrm{alg}}=2t$.
Tymczasem podział
\begin{equation}
  \{6t+1\},\qquad \{3t,3t\},\qquad \{2t,2t,2t\}
\end{equation}
ma sumy $(6t+1,6t,6t)$, czyli rozpiętość $1$.
Ponieważ suma wejścia $18t+1$ nie jest podzielna przez $3$,
rozpiętość $0$ jest niemożliwa i $\Delta^*=1$.
Stosunek $\Delta_{\mathrm{alg}}/\Delta^*=2t$ rośnie bez ograniczeń.

\subsubsection{Złożoność czasowa}

Analizujemy koszt bitowy, przyjmując najgorszy koszt $O(b)$
porównania dwóch wartości wejściowych i $O(s)$ porównania lub dodania
sum. Sortowanie przez scalanie par z zachowaniem ich oryginalnych
indeksów zajmuje $O(n\log n\,b)$ czasu; przy remisie dodatkowe
porównanie indeksów ma stały koszt w przyjętym modelu.
Inicjalizacja części kosztuje $O(kb)$, a budowa kopca $O(ks)$.

Po inicjalizacji pozostaje $n-k$ elementów. Każdy z nich wymaga
jednego dodawania na liczbach do $s$ bitów oraz operacji
\texttt{EXTRACT-MIN} i \texttt{INSERT}. Każda operacja kopca wykonuje
$O(\log k)$ porównań sum, co kosztuje $O(s\log k)$ czasu bitowego.
Wyznaczenie granic wszystkich bloków wymaga łącznie $O(n)$
operacji na indeksach. Otrzymujemy
\begin{equation}
  T(n,k,b)=O\bigl(n\log n\,b+ks+(n-k)s(1+\log k)\bigr),
  \qquad s=b+\lceil\log_2 n\rceil.
\end{equation}
Dla $2\leq k\leq n$ jest to w szczególności
\begin{equation}
  \boxed{T(n,k,b)=O\bigl(n\log n\,(b+\log n)\bigr).}
\end{equation}

\section{Kompilacja i uruchamianie programu}

Program kompilujemy poleceniem:
\begin{verbatim}
gcc -std=c17 main.c big_uint.c number_sort.c min_heap.c partition.c partition_exact.c -o taio
\end{verbatim}

Uruchomienie wymaga podania trybu algorytmu, pliku wejściowego
i liczby podzbiorów \(k\):
\begin{verbatim}
./taio A dane.txt 3
./taio D dane.txt 3
\end{verbatim}

Tryb \texttt{A} uruchamia algorytm aproksymacyjny,
a \texttt{D} algorytm dokładny. Wynik zostaje zapisany
w pliku \texttt{dane\_output.txt}.

\section{Testy obliczeniowe}

\subsection{Sposób przeprowadzenia i powtórzenia pomiarów}

Mierzymy czas działania oddawanego programu \texttt{taio} uruchamianego z~linii poleceń.
Dane tworzy generator \texttt{gen} (plik \texttt{gen.c}, czyste C17), który dla tych samych
argumentów zapisuje zawsze identyczny plik we wspólnym formacie wejściowym (własny generator
liczb pseudolosowych \emph{splitmix64} z~podanym ziarnem). Pomiary przeprowadza skrypt
\texttt{bench.py} (Python~3, wyłącznie biblioteka standardowa). Każdy punkt pomiarowy
\((n,k,\text{dane})\) jest wyznaczany tak samo:
\begin{enumerate}
    \item \texttt{gen} zapisuje plik z~\(R\) zbiorami po \(n\) liczb,
    \item \texttt{taio} przetwarza cały plik w~jednym uruchomieniu; mierzymy czas \(t\)
          uruchomienia, a gdy \(t<2\)\,s, pierwsze uruchomienie jest rozgrzewką i~bierzemy
          medianę z~trzech kolejnych,
    \item czas na jeden zbiór wynosi \(T=(t-T_0)/R\), gdzie \(T_0\) jest czasem startu
          programu (mediana z~21 uruchomień na pliku z~jednym zbiorem dwuelementowym);
          \(R\) zwiększamy, dopóki \(t\geq 0{,}5\)\,s, więc \(T\) jest \textbf{średnim czasem
          dla zbiorów tej samej wielkości}, a narzut startu procesu jest pomijalny,
    \item każdy plik wynikowy jest sprawdzany: podzbiory pokrywają zbiór, sumy i~porządek
          wierszy są poprawne, dla algorytmu aproksymacyjnego \(\Delta\leq a_{\max}\).
\end{enumerate}
Czas \(T\) obejmuje wczytanie zbioru, obliczenie i~zapis wyniku. Seria dla algorytmu
dokładnego kończy się na pierwszym \(n\), dla którego \(T>10\)\,s (w~serii z~kilkoma
porządkami danych -- dla któregokolwiek z~nich). Stosujemy trzy rodziny
danych: \texttt{pes} -- rodzina pesymistyczna z~analizy algorytmu dokładnego
(\(a_1=\cdots=a_{n-1}=k\), \(a_n=k(n+1)\), podział idealny nie istnieje),
\texttt{los} -- liczby losowe z~\([1,2^b-1]\) z~sumą podzielną przez \(k\) (ostatnia liczba
korygowana o~mniej niż \(k\)), \texttt{idealny} -- \(k\) grup o~równych sumach wymieszanych
losowo (podział idealny istnieje, \(\Delta^*=0\)). Dane \texttt{los} i~\texttt{idealny}
mogą być zapisane w~kolejności losowej, nierosnącej lub niemalejącej; dla tego samego ziarna
jest to ten sam multizbiór.

Pomiary powtarzamy poleceniami (Windows, PowerShell, w~katalogu ze źródłami; w~systemie
Linux odpowiednio bez rozszerzenia \texttt{.exe} i~z~\texttt{python3} zamiast \texttt{py}):
{\footnotesize
\begin{verbatim}
gcc -std=c17 -O2 main.c big_uint.c number_sort.c min_heap.c partition.c partition_exact.c -o taio.exe
gcc -std=c17 -O2 gen.c big_uint.c -o gen.exe
py bench.py
\end{verbatim}}
\noindent Pojedynczą serię można powtórzyć poleceniem \texttt{py bench.py --serie T2}.
Wyniki trafiają do plików \texttt{wyniki/T1.csv}, \dots, \texttt{wyniki/T5.csv}
(każdy wiersz zawiera m.in.\ parametry generatora z~ziarnem, \(R\), \(t\), \(T_0\) i~\(T\)).
Wykresy w~tym dokumencie są rysowane bezpośrednio z~tych plików (pakiet \texttt{pgfplots}),
więc ponowna kompilacja dokumentu po pomiarach odświeża wykresy. Przedstawione wyniki
uzyskano na komputerze z~procesorem AMD Ryzen 7 9700X (Windows~11, pomiary w~środowisku
WSL2 z~Ubuntu 24.04, gcc 13.3, Python 3.12, \(T_0\approx0{,}6\)\,ms). Na innym komputerze
czasy bezwzględne będą inne, ale kształt zależności pozostaje ten sam.

\pgfplotstableread[col sep=comma]{wyniki/T1.csv}\tabTjeden
\pgfplotstableread[col sep=comma]{wyniki/T2.csv}\tabTdwa
\pgfplotstableread[col sep=comma]{wyniki/T3.csv}\tabTtrzy

\subsection{Algorytm dokładny}

Na danych \texttt{pes} algorytm zawsze przegląda wszystkie \(k^n\) kandydatów
(rys.~\ref{fig:t1}). W~skali logarytmicznej punkty układają się w~proste: zwiększenie \(n\)
o~jeden wydłuża czas średnio \(2{,}05\), \(3{,}13\), \(4{,}29\) i~\(5{,}31\) razy dla
\(k=2,3,4,5\), czyli nieco bardziej niż \(k\) razy, bo rośnie też liczba operacji na
jednego kandydata. Prawy wykres pokazuje czas przypadający na jednego kandydata,
\(T/k^n\): rośnie on w~przybliżeniu liniowo z~\(n\) (od około 40 do 90\,ns), co odpowiada
kosztowi \(\Theta(n+k)\) sprawdzenia kandydata i~potwierdza oszacowanie
\(\Theta\bigl((n+k)\,k^n\bigr)\). Na składnik niezależny od \(n\) składa się m.in.\ obsługa
\(k\) sum jako liczb dowolnej precyzji, których pamięć w~implementacji jest zwalniana
i~przydzielana na nowo dla każdego kandydata. Największe \(n\), dla których czas nie
przekracza 10\,s, wynoszą 26, 17, 13 i~11 odpowiednio dla \(k=2,3,4,5\).

\begin{figure}[htbp]
    \centering
    \begin{tikzpicture}
        \begin{semilogyaxis}[dwa, name=lewy, xlabel={liczba elementów \(n\)},
                ylabel={czas na zbiór \(T\) [s]}, legend pos=north west]
            \addplot+[discard if not={k}{2}] table[x=n, y=t_na_zbior_s]{\tabTjeden};
            \addlegendentry{\(k=2\)}
            \addplot+[discard if not={k}{3}] table[x=n, y=t_na_zbior_s]{\tabTjeden};
            \addlegendentry{\(k=3\)}
            \addplot+[discard if not={k}{4}] table[x=n, y=t_na_zbior_s]{\tabTjeden};
            \addlegendentry{\(k=4\)}
            \addplot+[discard if not={k}{5}] table[x=n, y=t_na_zbior_s]{\tabTjeden};
            \addlegendentry{\(k=5\)}
        \end{semilogyaxis}
        \begin{axis}[dwa, prawy, xlabel={liczba elementów \(n\)},
                ylabel={\(T/k^n\) [ns]}, ymin=0]
            \addplot+[discard if not={k}{2}]
                table[x=n, y expr=\thisrow{t_na_zbior_s}*1e9/(\thisrow{k}^\thisrow{n})]{\tabTjeden};
            \addplot+[discard if not={k}{3}]
                table[x=n, y expr=\thisrow{t_na_zbior_s}*1e9/(\thisrow{k}^\thisrow{n})]{\tabTjeden};
            \addplot+[discard if not={k}{4}]
                table[x=n, y expr=\thisrow{t_na_zbior_s}*1e9/(\thisrow{k}^\thisrow{n})]{\tabTjeden};
            \addplot+[discard if not={k}{5}]
                table[x=n, y expr=\thisrow{t_na_zbior_s}*1e9/(\thisrow{k}^\thisrow{n})]{\tabTjeden};
        \end{axis}
    \end{tikzpicture}
    \caption{Algorytm dokładny na danych \texttt{pes} (pełny przegląd). Po lewej czas na
             zbiór w~skali logarytmicznej, po prawej czas przypadający na jednego kandydata.}
    \label{fig:t1}
\end{figure}

Na danych losowych \texttt{los} (\(b=20\), czyli wartości do około \(10^6\)) czas zależy
przede wszystkim od tego, czy podział idealny istnieje (rys.~\ref{fig:t2}). Dla \(k=3\)
żaden z~badanych zbiorów nie miał podziału idealnego, więc we wszystkich porządkach czas
jest równy czasowi pełnego przeglądu (linia \texttt{pes}). Dla \(k=2\) do \(n=18\) podział
idealny ma mniej niż 2\% zbiorów, w~przedziale \(n=19\)--\(24\) odsetek ten szybko rośnie
(w~próbach po 10--20 zbiorów od 0\% do 100\%), a od \(n=25\) podział idealny mają wszystkie
zbiory. Algorytm kończy się wtedy na pierwszym znalezionym podziale i~czas rośnie
znacznie wolniej niż \(2^n\), co pozwoliło dojść do \(n=33\) (pełny przegląd trwałby dla
tego \(n\) około 15\,min). O~czasie decyduje wówczas kolejność elementów. Dla \(n=33\)
średni czas wynosi 10,1\,s przy porządku nierosnącym (takim, jakiego używa algorytm
aproksymacyjny), 1,0\,s przy losowym i~0,48\,s przy niemalejącym. Wynika to z~kolejności
przeglądania kandydatów: najszybciej zmienia się przydział ostatnich elementów, więc gdy na
końcu stoją największe liczby, podział o~równych sumach zostaje znaleziony wcześniej.
Dla największych \(n\) każdy punkt jest średnią z~10 zbiorów, stąd widoczny rozrzut.

\begin{figure}[htbp]
    \centering
    \begin{tikzpicture}
        \begin{semilogyaxis}[dwa, name=lewy, title={\(k=2\)},
                xlabel={\(n\)}, ylabel={czas na zbiór \(T\) [s]}, legend pos=north west]
            \addplot+[discard if not={k}{2}, discard if not={porzadek}{los}]
                table[x=n, y=t_na_zbior_s]{\tabTdwa};
            \addlegendentry{losowy}
            \addplot+[discard if not={k}{2}, discard if not={porzadek}{nier}]
                table[x=n, y=t_na_zbior_s]{\tabTdwa};
            \addlegendentry{nierosnący}
            \addplot+[discard if not={k}{2}, discard if not={porzadek}{niem}]
                table[x=n, y=t_na_zbior_s]{\tabTdwa};
            \addlegendentry{niemalejący}
            \addplot[gray, dashed, discard if not={k}{2}]
                table[x=n, y=t_na_zbior_s]{\tabTjeden};
            \addlegendentry{\texttt{pes}}
        \end{semilogyaxis}
        \begin{semilogyaxis}[dwa, prawy, title={\(k=3\)}, xlabel={\(n\)},
                ylabel={\(T\) [s]}, legend pos=north west]
            \addplot+[discard if not={k}{3}, discard if not={porzadek}{los}]
                table[x=n, y=t_na_zbior_s]{\tabTdwa};
            \addlegendentry{losowy}
            \addplot+[discard if not={k}{3}, discard if not={porzadek}{nier}]
                table[x=n, y=t_na_zbior_s]{\tabTdwa};
            \addlegendentry{nierosnący}
            \addplot+[discard if not={k}{3}, discard if not={porzadek}{niem}]
                table[x=n, y=t_na_zbior_s]{\tabTdwa};
            \addlegendentry{niemalejący}
            \addplot[gray, dashed, discard if not={k}{3}]
                table[x=n, y=t_na_zbior_s]{\tabTjeden};
            \addlegendentry{\texttt{pes}}
        \end{semilogyaxis}
    \end{tikzpicture}
    \caption{Algorytm dokładny na danych \texttt{los} (\(b=20\)) w~trzech porządkach
             wejścia; linia przerywana -- dane \texttt{pes} z~rys.~\ref{fig:t1}. Dla \(k=3\)
             wszystkie krzywe pokrywają się.}
    \label{fig:t2}
\end{figure}

\subsection{Algorytm aproksymacyjny}

W~wymaganym zakresie \(5\leq n\leq 200\) (rys.~\ref{fig:t3}, po lewej) czas na zbiór
rośnie praktycznie liniowo, od około 1,6\,\(\mu\)s dla \(n=5\) do 54--59\,\(\mu\)s dla
\(n=200\). Wpływ \(k\) jest mały: dla \(k=10\) czas jest zwykle o~kilka do kilkunastu
procent dłuższy niż dla \(k=2\). Dane już posortowane nierosnąco są przetwarzane zwykle
o~10--15\% szybciej. Prawy wykres obejmuje zakres do
\(n=10^6\); linia przerywana to \(c\,n\log_2 n\) ze stałą \(c=20{,}65\)\,ns dobraną tak,
aby przechodziła przez punkt \(n=10^5\). Tabela~\ref{tab:t3} podaje iloraz
\(T/(n\log_2 n)\) dla wybranych \(n\).

\begin{table}[htbp]
    \centering
    \begin{tabular}{rrrrr}
        \toprule
        & \multicolumn{2}{c}{\(T\) [ms]} & \multicolumn{2}{c}{\(T/(n\log_2 n)\) [ns]} \\
        \(n\) & dane losowe & posortowane & dane losowe & posortowane \\
        \midrule
        \(10^3\) & 0,284 & 0,227 & 28,5 & 22,8 \\
        \(10^4\) & 3,03 & 2,27 & 22,8 & 17,1 \\
        \(10^5\) & 34,3 & 23,5 & 20,7 & 14,1 \\
        \(10^6\) & 664 & 271 & 33,3 & 13,6 \\
        \bottomrule
    \end{tabular}
    \caption{Algorytm aproksymacyjny, \(k=3\), \(b=64\): czas na zbiór dla dużych \(n\).}
    \label{tab:t3}
\end{table}

Iloraz \(T/(n\log_2 n)\) maleje aż do \(n\approx2\cdot10^5\), bo w~mierzonym czasie
przeważa składnik liniowy: wczytanie i~zapis liczb w~postaci dziesiętnej oraz przydział
pamięci dla każdej liczby. Sam składnik \(n\log n\) (sortowanie) jest więc w~tym zakresie
słabo widoczny. Dla danych losowych przy \(n\geq5\cdot10^5\) iloraz rośnie, a dla danych
posortowanych pozostaje stały (około 13,5\,ns). Przyczyną są najpewniej chybienia w~pamięci
podręcznej podczas sortowania losowo ułożonych liczb, a nie sam algorytm. Pomiary są
zgodne z~oszacowaniem \(O\bigl(n\log n\,(b+\log n)\bigr)\), a zbiór z~milionem liczb
algorytm przetwarza w~czasie poniżej 1\,s.

\begin{figure}[htbp]
    \centering
    \begin{tikzpicture}
        \begin{axis}[dwa, name=lewy, xlabel={\(n\)}, ylabel={czas na zbiór \(T\) [ms]},
                restrict x to domain=0:200, xmin=0, xmax=205, ymin=0, legend pos=north west]
            \addplot+[mark size=1pt, discard if not={k}{2}, discard if not={porzadek}{los}]
                table[x=n, y expr=\thisrow{t_na_zbior_s}*1000]{\tabTtrzy};
            \addlegendentry{\(k=2\)}
            \addplot+[mark size=1pt, discard if not={k}{3}, discard if not={porzadek}{los}]
                table[x=n, y expr=\thisrow{t_na_zbior_s}*1000]{\tabTtrzy};
            \addlegendentry{\(k=3\)}
            \addplot+[mark size=1pt, discard if not={k}{10}, discard if not={porzadek}{los}]
                table[x=n, y expr=\thisrow{t_na_zbior_s}*1000]{\tabTtrzy};
            \addlegendentry{\(k=10\)}
            \addplot+[mark size=1pt, discard if not={k}{3}, discard if not={porzadek}{nier}]
                table[x=n, y expr=\thisrow{t_na_zbior_s}*1000]{\tabTtrzy};
            \addlegendentry{\(k=3\), posortowane}
        \end{axis}
        \begin{loglogaxis}[dwa, prawy, xlabel={\(n\)}, ylabel={\(T\) [s]},
                legend pos=north west]
            \addplot+[mark size=1pt, discard if not={k}{3}, discard if not={porzadek}{los}]
                table[x=n, y=t_na_zbior_s]{\tabTtrzy};
            \addlegendentry{\(k=3\)}
            \addplot+[mark size=1pt, discard if not={k}{3}, discard if not={porzadek}{nier}]
                table[x=n, y=t_na_zbior_s]{\tabTtrzy};
            \addlegendentry{\(k=3\), posortowane}
            % n = 10^x, c n log2 n = 20,65 ns * 10^x * x * log2(10)
            \addplot[black, dashed, domain=0.7:6, samples=30]
                ({10^x}, {2.065e-8*3.321928*x*10^x});
            \addlegendentry{\(c\,n\log_2 n\)}
        \end{loglogaxis}
    \end{tikzpicture}
    \caption{Algorytm aproksymacyjny na danych \texttt{los} (\(b=64\)): po lewej zakres
             \(5\leq n\leq 200\), po prawej \(5\leq n\leq 10^6\) w~skali log-log.}
    \label{fig:t3}
\end{figure}

Tabela~\ref{tab:t4} pokazuje wpływ długości liczb \(b\) przy \(n=10^5\) i~\(k=3\).
Do \(b=256\) czas rośnie wolniej niż liniowo, bo przeważają koszty niezależne od \(b\),
a powyżej szybciej: podwojenie \(b\) z~512 do 1024 wydłuża czas 2,6 razy. Wynika to
z~zamiany zapisu dziesiętnego na binarny i~z~powrotem przy wczytywaniu i~zapisie. Koszt
tej zamiany dla jednej liczby rośnie jak \(b^2\), podczas gdy operacje samego algorytmu
(porównania i~dodawania sum) mają koszt \(O(b)\).

\begin{table}[htbp]
    \centering
    \begin{tabular}{lrrrrrrr}
        \toprule
        \(b\) [bity] & 16 & 32 & 64 & 128 & 256 & 512 & 1024 \\
        \midrule
        \(T\) [ms] & 23,3 & 27,5 & 35,0 & 48,6 & 84,6 & 187,5 & 494,6 \\
        \bottomrule
    \end{tabular}
    \caption{Algorytm aproksymacyjny, \(n=10^5\), \(k=3\), dane \texttt{los}: czas na zbiór
             w~zależności od długości liczb (\texttt{wyniki/T4.csv}).}
    \label{tab:t4}
\end{table}

\subsection{Jakość rozwiązań algorytmu aproksymacyjnego}

Na danych \texttt{idealny} (\(b=30\), po 100 zbiorów dla każdego \(n\)) podział idealny
istnieje z~konstrukcji, czyli \(\Delta^*=0\). Algorytm aproksymacyjny znajduje go jednak
tylko w~6\% (\(k=2\)) i~5\% (\(k=3\)) zbiorów dla \(n=10\), a dla \(n\geq20\) w~żadnym
(tab.~\ref{tab:t5}). Średni iloraz \(\Delta_{\mathrm{alg}}/a_{\max}\) jest natomiast mały
i~maleje w~przybliżeniu jak \(1/n\), więc gwarancja \(\Delta_{\mathrm{alg}}\leq a_{\max}\)
jest w~praktyce spełniona z~dużym zapasem. Z~kolei algorytm dokładny zawsze znajduje
podział idealny, jeśli on istnieje, ale tylko dla niewielkich \(n\).

\begin{table}[htbp]
    \centering
    \begin{tabular}{lrrrrr}
        \toprule
        \(n\) & 10 & 20 & 50 & 100 & 200 \\
        \midrule
        \(k=2\): średnie \(\Delta_{\mathrm{alg}}/a_{\max}\) & 0,080 & 0,042 & 0,015 & 0,0069 & 0,0034 \\
        \(k=2\): zbiory z~\(\Delta_{\mathrm{alg}}=0\) & 6\% & 0\% & 0\% & 0\% & 0\% \\
        \(k=3\): średnie \(\Delta_{\mathrm{alg}}/a_{\max}\) & 0,127 & 0,084 & 0,029 & 0,013 & 0,0073 \\
        \(k=3\): zbiory z~\(\Delta_{\mathrm{alg}}=0\) & 5\% & 0\% & 0\% & 0\% & 0\% \\
        \bottomrule
    \end{tabular}
    \caption{Jakość algorytmu aproksymacyjnego na danych z~podziałem idealnym
             (\(\Delta^*=0\), \texttt{wyniki/T5.csv}).}
    \label{tab:t5}
\end{table}


\section{TODO: Wnioski}

\section{TODO: podział pracy}

\end{document}